\chapter{Método, fuentes y límites}

\label{cap:metodo}
\section{Los tres registros}

El libro distingue tipográficamente tres niveles y no los mezcla. \textbf{Documento}: sentencias, ordenanzas, decretos, resoluciones, boletines oficiales, expedientes administrativos, publicaciones institucionales. Se cita con identificación completa. \textbf{Prensa}: cobertura periodística verificable, con medio y fecha, consignada siempre como afirmación del medio y no como hecho establecido. \textbf{Inferencia del autor}: va en caja gris, nunca fuera de ella.

\section{El corpus}

\textbf{Judicial.} La causa \textit{Lazarte Vigabriel y otros c/ Municipalidad de La Caldera; Ministerio de la Producción, Trabajo y Desarrollo Sustentable --- Amparo Colectivo} (CUIJ 17-00023860-0; piezas pertenecientes INC 23860; incidente INC 23860/1; ejecución, Expte.\ 780280), ante el Juzgado de Minas y en lo Comercial de Registro del Distrito Judicial Centro, es la columna vertebral del libro.

\textbf{Sobre la numeración del expediente principal hay una discrepancia que este libro no resuelve.} Una providencia de agosto de 2022 lo identifica como \textbf{MIN 23860/19}; otras piezas del corpus lo consignan como \textbf{EMI 23860/2020}. Ambas formas se conservan tal como aparecen en sus fuentes. Esa misma providencia deja constancia de que en 2022 el expediente principal \textbf{se encontraba radicado ante la Corte de Justicia de la Provincia}, razón por la cual la etapa de cumplimiento se tramitó en piezas separadas.

Conviene anotar que \textbf{la sentencia la dicta el Juzgado de Minas} y que la etapa de cumplimiento tramita en piezas separadas, con carátula de amparo colectivo y numeración propia (INC 23860/1 y 780280); el capítulo \ref{cap:amparo} explica esa relación. Su contenido se reconstruye aquí a partir de la ficha del caso elaborada por el Observatorio Argentino de Litigación Climática de la Universidad Nacional del Litoral, de los partes de prensa del Poder Judicial de Salta y de la consulta al sistema de expediente digital; el texto íntegro de la sentencia no fue consultado, y ésa es la principal debilidad documental de este libro. \textbf{La única transcripción disponible es parcial y la eligió el propio obligado}: el informe que la Provincia presentó en 2022 cita «la parte pertinente» de los mandatos que debía cumplir --- del punto VII sólo el inciso a) y sólo su primer párrafo ---, y su estructura sugiere que hay al menos cuatro mandatos más cuyo texto no consta. Se le suman el amparo por la cuenca del río Mojotoro de 2019, el amparo rechazado contra la Municipalidad de Vaqueros por el loteo \textit{Lares de la Inmaculada}, y como caso comparado el peritaje sobre el barrio privado \textit{Los Maitines} en San Lorenzo.

\textbf{Normativo.} Se obtuvo y leyó el texto operativo de: el Código de Ordenamiento Territorial de La Caldera (152 artículos, ocho anexos, atlas de ocho planos); el Informe Final y el Proyecto de Ordenanza del Plan de Desarrollo Urbano Ambiental que ese Código reglamenta; el Código de Planeamiento Urbano Ambiental de Vaqueros en su texto de 2026; la Ley 2616 de 1951; el Decreto 399/20 de creación del Plan ``Mi Lote''; las Resoluciones 82/16 y 86/16 del Ministerio y 7/16, 8/16 y 9/16 de la Secretaría de Tierra y Bienes, con los anexos de la 82/16, la 7/16, la 8/16 y la 9/16 ---copias certificadas que el repositorio del Boletín ofrece sólo como imagen---; el Decreto 4.414/08 que encabeza esa cadena; los decretos de designación 120/15 y 64/23; el Decreto 844/21 de traslado orgánico; el Decreto 63/25 de adjudicación y su anexo nominal.

\textbf{Cómo se transcribe.} Las citas de documentos reproducen el texto del original, con dos
convenciones declaradas. \textbf{Se moderniza la acentuación} ---el Boletín histórico
imprime «Nacion», «rio», «hectàreas», «areas»; los ejemplos son de 1908 a 1912, y la
convención rige para todo el corpus anterior a 1947--- y \textbf{se unen las palabras partidas por el
corte de renglón}, porque ni una cosa ni la otra pertenecen al texto sino a la imprenta. En
cambio \textbf{las erratas de palabra se conservan tal cual, con \textit{[sic]}}: así van «el
Augosto de Arias» junto a «el Angosto de Arrieta» en el mismo decreto de 1912, «Rejis» por
«Regis», «Depto. de la Caldera» y las dos superficies distintas de una misma finca en un mismo
aviso. Donde el original se contradice, se transcriben las dos formas y no se elige.

\textbf{Histórico.} El archivo digital del Boletín Oficial de la Provincia, desde su número 1 del 15 de octubre de 1908. Se leyeron \textbf{edición por edición} treinta y seis tramos, que el apéndice \ref{cap:fuentes} detalla uno por uno con sus ausencias de edición y de hoja: las 196 primeras ediciones, de octubre de 1908 a octubre de 1910; los años \textbf{1911 a 1921} completos, del número 220 al 888; el \textbf{primer semestre de 1922}, del 890 al 916, que es hasta donde llega el repositorio ---de la edición 917 en adelante no hay PDF---; los años \textbf{1923 a 1932} completos, del 939 al 1.461; las ediciones que van del número 962 al 2.464, de junio de 1923 a diciembre de 1945, de las que faltan treinta y una; y, al cierre, los años \textbf{1937 a 1946} completos, hasta el número 2.744, y \textbf{1947 a 1950}, de la edición 2.869 en adelante, estos cuatro con barrido automático y \textbf{sin lectura sobre la imagen}, según declara el apéndice \ref{cap:fuentes}. Se barrió además por término de búsqueda el período 2013--2015, y de manera puntual el resto. El apéndice \ref{cap:fuentes} detalla el método. De ese corpus provienen los partidos policiales de 1909 y 1911, los decretos de defensa contra el río de 1910 y 1911, las rendiciones municipales de 1908 y 1922, los pedimentos mineros de 1908 y 1925, los deslindes de 1910 y 1921, la concesión de agua de 1924, la fundición de Mojotoro y el racionamiento del Wierna de 1946.

\textbf{Hemerográfico.} La cronología del apéndice \ref{cap:cronologia} reúne el barrido 2015--2026. Los medios con cobertura sostenida del proceso son \textit{Página/12} en su edición salteña, \textit{El Tribuno}, \textit{Cuarto Poder}, \textit{Que Pasa Salta}, \textit{Cuarto}, \textit{Nuevo Diario de Salta} y \textit{Salta 4400}.

\textbf{Académico.} Literatura sobre Vaqueros como caso de transformación periurbana; trabajo sobre urbanización de la pobreza en el aglomerado Gran Salta; y un relevamiento de la universidad confesional provincial que identificó, para Salta capital y municipios vecinos, cinco urbanizaciones privadas y once clubes de campo aprobados por la Dirección General de Inmuebles, frente a cincuenta y cuatro asentamientos.

\begin{cautela}
Ese contraste --- dieciséis urbanizaciones cerradas contra cincuenta y cuatro asentamientos en el mismo aglomerado --- es el marco de todo el libro. No hay un mercado de suelo que crece: hay dos mercados que crecen en paralelo, con regímenes jurídicos, ritmos de provisión de servicios y capacidades de gestión radicalmente distintos, sobre el mismo aglomerado.
\end{cautela}

\section{Un sesgo de la fuente que terminó siendo un hallazgo}

El obstáculo apareció primero en cinco instrumentos: el Reglamento de Fraccionamiento de Tierras de Vaqueros (Ordenanza 437), la nómina de adjudicatarios del Decreto 63/25, el baremo de puntaje de la Resolución 8/16, los procedimientos de sorteo de la Resolución 9/16 y el Reglamento de Asignación de la Resolución 82/16.

Al ampliarse el relevamiento la lista creció. Se sumaron el contrato de promoción turística que ratifica el Decreto 3069/14, con el detalle de las exenciones concedidas; las coordenadas de dos determinaciones de línea de ribera, remitidas a la sede del organismo; y el informe técnico que funda la contratación por vía abreviada del Plan de Manejo.

Con los que se agregaron después ---el anexo del Decreto 1682/19, el de la Resolución 023/14, los de la Resolución 567D/26 y el convenio del Decreto 1433/15--- la lista llegó a trece. De los trece, \textbf{ocho resultaron estar publicados} como archivos vinculados en la edición web del Boletín ---la nómina del Decreto 63/25, el reglamento de la Resolución 82/16, los anexos de las Resoluciones 8/16 y 9/16 y los cuatro que se sumaron al final---, varios de ellos sólo como imagen escaneada; los \textbf{cinco restantes} siguen fuera, contando por separado las dos determinaciones de línea de ribera. El capítulo \ref{cap:opacidad} da la lista definitiva ---seis instrumentos, de los que al final del trabajo se obtuvo uno, el reglamento de Vaqueros---, que son esos cinco ---la Ordenanza 437, las Resoluciones 142/14 y 383/14, el Decreto 3069/14 y el informe técnico del Plan de Manejo--- más el Decreto 463/15, hallado al verificarlos.

Los anexos no son secretos: existen, integran los actos y son accesibles por pedido. Pero contienen, sin excepción, el contenido decisorio --- qué obras debe hacer un loteador, cuántos puntos vale una discapacidad, cuántos días tiene un vecino para impugnar un padrón --- mientras la parte publicada contiene la envoltura. Quien insiste los consigue; quien lee el Boletín Oficial de buena fe como texto, no: varios de ellos sólo existen allí como imagen.

Este apartado se escribió, en una primera versión, como una advertencia práctica para quien continuara la investigación. La repetición del obstáculo lo convirtió en objeto: el capítulo \ref{cap:opacidad} lo retoma como segunda tesis del libro.

\begin{cautela}
Esto tiene una consecuencia metodológica que el libro debió aprender por error. Trabajar sobre una norma sin su anexo produce la impresión de un vacío donde hay una regla. Este relevamiento afirmó, en una etapa intermedia, que la asignación directa de lotes fiscales era una excepción admitida sin reglamentar; al obtener el anexo faltante resultó ser una de las cláusulas mejor acotadas del régimen, con destinatario definido, informe fundado, doble firma, período de tachas y tope del dos por ciento.

La regla que se sigue de aquí en adelante, y que se recomienda a quien continúe: ninguna afirmación sobre ausencia de norma sin haber agotado la cadena de remisiones.
\end{cautela}

\section{Una trampa de homonimia}

Durante el relevamiento apareció un sitio de noticias con dominio \texttt{lacaldera.com.ar}, con secciones de política y judiciales, que publica notas sobre el Partido Justicialista, denuncias de aprietes a intendentes y licitaciones sospechosas, y una nota de junio de 2025 sobre un dirigente de apellido Solanas. Nada de eso pertenece a este objeto: es un medio de Entre Ríos, con cobertura de Paraná, Concordia y Diamante, y la coincidencia con el nombre del loteo caldereño es puramente léxica.

\begin{cautela}
Se consigna porque describe con exactitud cómo se corrompe una investigación de este tipo. Un dominio que coincide con el nombre del objeto, un apellido que coincide con el nombre de un emprendimiento y un registro retórico perfectamente compatible con la tesis. Tres coincidencias formales y ninguna referencial.

De ahí se derivan dos reglas que el libro aplica sin excepción. Ninguna fuente entra al corpus sin verificación de jurisdicción. Y ningún cotejo de apellidos sustituye a una identificación documental: los apellidos involucrados en este territorio --- Quipildor, Chocobar, Mamaní, Duarte, Tolaba --- están entre los más frecuentes de la provincia, y cruzarlos contra una imputación produce coincidencia, no prueba.
\end{cautela}

Y una tercera, de protección de datos. Los instrumentos que el libro cita consignan, con frecuencia, el número de documento de las personas físicas que nombran. \textbf{Este libro no los reproduce.} Cuando la identificación de una persona depende de ese número ---que el gerente de una sociedad sea el mismo que figura en un acta electoral, o que dos documentos sean correlativos---, se dice que el cotejo se hizo y cuál fue su resultado, sin transcribir el número.

El mismo criterio rige para los nombres. Cuando un registro nominal identifica a personas
que pueden estar vivas y que no ejercen función pública ni son parte de un acto que este
libro examina en su calidad de tales, no se reproducen sus nombres: se dice qué contiene la
lista, qué se hizo con ella y qué se leyó en agregado. Así se procede con el anexo nominal
del Decreto 63/25, que identifica con nombre y documento a un centenar de adjudicatarios de
tierra fiscal (capítulo \ref{cap:tierrafiscal}), y con la partida del revalúo de 1944 cuyos
titulares eran menores de edad, que hoy podrían estar vivos (capítulo \ref{cap:fincas}). A
los titulares adultos de ese mismo revalúo, en cambio, el libro los nombra: son los
propietarios de un padrón público de hace más de ochenta años, y esa nómina es el objeto del
capítulo.

\section{Lo que este libro no puede hacer}

No tiene acceso registral: sin consulta a la Dirección General de Inmuebles, la matriz de titularidades del apéndice \ref{cap:matriz} queda incompleta, y sus casilleros vacíos se publican vacíos. No tiene el digesto municipal de La Caldera, que no publica sus ordenanzas de manera accesible; el de Vaqueros, en cambio, está en línea desde al menos febrero de 2024, y de él salió al cierre el anexo de la Ordenanza 437 (capítulo \ref{cap:vaqueros}). El sitio oficial de La Caldera estuvo fuera de servicio durante buena parte del relevamiento y volvió a estarlo después; sus partes de prensa de febrero de 2026, que este libro sí utilizó, se consultaron en una de las ventanas en que estuvo disponible. No tiene trabajo de campo: no hay entrevistas propias, sólo dos conversaciones informales con vecinos, de septiembre de 2026, que el capítulo \ref{cap:poblacion} registra como testimonio y no como dato. \textbf{Sí existen, en cambio, las que hizo otro}: el Plan de Manejo incorporado al expediente judicial se apoya en encuestas estructuradas y semiestructuradas a vecinos del departamento, realizadas entre febrero y julio de 2024 y ofrecidas bajo anonimato porque los encuestados manifestaron temor a represalias. Los formularios quedaron en resguardo de la consultora y a disposición del juzgado. El capítulo \ref{cap:amparo} lo documenta. Y existe además un libro escrito por habitantes del departamento: \textit{Nuestro tiempo redondo}, publicado en 2019 por la Comunidad Kolla Kondorwaira y editado por el municipio, que el capítulo \ref{cap:resistencias} incorpora. \textbf{Ninguna de las dos cosas reemplaza el trabajo de campo que este libro no hizo}, pero ambas muestran que el material existe. Y tiene un hueco temporal grande y declarado: la lectura edición por edición termina en diciembre de 1950 ---1947 a 1950 se barrieron al cierre, con la reserva que el apéndice \ref{cap:fuentes} declara: pre-informe automático y sin lectura sobre la imagen--- y el barrido sistemático siguiente empieza en 2013, de modo que \textbf{los sesenta y dos años que van de 1951 a 2012 son un período que este trabajo no leyó de corrido}. Es el período en que se crea el municipio de Vaqueros, se expropia el Campo Alegre, se dicta el régimen de loteos de 1973, se resuelve la sucesión Serrey\index{Serrey, Carlos} y ocurre el boom de los noventa. Hay además un segundo hueco, más corto y anterior, y al cierre quedó reducido a la mitad: el \textbf{primer semestre de 1922} se leyó edición por edición y 1923 se leyó entero, de modo que lo que queda sin lectura continua es el \textbf{segundo semestre de 1922}, del que el repositorio no tiene un solo PDF. El capítulo \ref{cap:planteo} ubica en el primero de los dos huecos, el de 1951 a 2012 ---con 1947 a 1950 barridos pero no leídos sobre la imagen---, la respuesta a cómo el municipio perdió la voz que tenía sobre el agua de su jurisdicción. Y no pudo leer en bloque la última década del Boletín Oficial: el repositorio descargable del organismo cubre de 1908 a 2015 y se interrumpe ahí, de modo que los años posteriores se consultaron acto por acto, sabiendo de antemano qué buscar. El capítulo \ref{cap:opacidad} desarrolla lo que eso significa.

Y hay un límite técnico que conviene declarar, porque afectó a este trabajo y afectará a quien lo continúe. \textbf{El reconocimiento óptico de caracteres colapsa las tablas escaneadas}: un topónimo, un nombre o una cifra que sólo existan dentro de una tabla son invisibles para cualquier búsqueda por texto. Este relevamiento dio por ausente el parámetro Yacones en un documento de más de doscientas páginas donde figura dos veces, las dos en tablas. \textbf{Antes de afirmar que un término no está en una fuente escaneada hay que mirar sus tablas con los ojos.}

Existe además un \textbf{repositorio digital de anexos} separado del boletín, que este trabajo consultó. Contiene material en volumen --- trescientas veintiséis páginas sólo en octubre de 2017 --- pero sus archivos llevan extensión de PDF sin serlo, algunos vienen sin capa de texto y otros sin marcas de página, de modo que citarlos exige reconocimiento óptico previo y no siempre permite dar referencia de página. Sobre las comunidades originarias, este trabajo corrigió su propio límite: sostenía no haber verificado su existencia, y el expediente 0010007-67527/2019-28 de Fiscalía de Estado prueba que \textbf{hay dos comunidades registradas en el departamento y que la zona alta de la cuenca es territorio de ocupación tradicional relevado por el INAI desde 2011}. El capítulo \ref{cap:resistencias} lo desarrolla. Lo que sigue sin verificarse es la existencia de familias con posesión veinteñal no comunitaria; no debe asumirse su ausencia.
